\AtBeginSection[]
{
    \begin{frame}
        \frametitle{Visão Geral}
        \tableofcontents[currentsection]
    \end{frame}
}

\usefonttheme{professionalfonts} % using non standard fonts for beamer
\usefonttheme{serif} % default family is serif
\usepackage{fontspec}
\setmainfont{XCharter}[
    UprightFont    = *-Roman,
    BoldFont       = *-Bold,
    ItalicFont     = *-Italic,
    BoldItalicFont = *-BoldItalic,
]

\usepackage{hyperref}
\usepackage{graphicx} % Allows including images
\usepackage{booktabs} % Allows the use of \toprule, \midrule and \bottomrule in tables
\input{./slides/SimplePlusTheme/beamerthemeSimplePlus.sty}
\input{./slides/SimplePlusTheme/beamerinnerthemeSimplePlus.sty}
\input{./slides/SimplePlusTheme/beamerfontthemeSimplePlus.sty}
\input{./slides/SimplePlusTheme/beamercolorthemeSimplePlus.sty}

\usepackage{listings}
\setbeamercovered{transparent}

\usepackage{biblatex}
\addbibresource{slides/reference.bib}
\renewcommand*{\bibfont}{\footnotesize}
% Tamanho dos slides de "Leituras Recomendadas": maior que a bibliografia
% completa (\bibfont acima, que rende \tiny via a redefinição de \footnotesize)
% para destacar as leituras curadas.
\newcommand{\leiturasize}{\small}

\usepackage[brazilian ]{babel}
\usepackage{csquotes}
\usepackage{pgfplots}
\pgfplotsset{compat=1.18}

\usepackage{emoji}
\usepackage{amsmath}
\usepackage[xcharter,bigdelims,vvarbb]{newtxmath}
% newtxmath's operators font requests the fontspec family `XCharter(0)' in OT1;
% without these substitutions math digits and operator names silently fall back
% to Computer Modern (LaTeX Font Warning: Font shape `OT1/XCharter(0)/m/n' undefined).
\DeclareFontFamily{OT1}{XCharter(0)}{}
\DeclareFontShape{OT1}{XCharter(0)}{m}{n}{<->ssub * XCharter-TLF/m/n}{}
\DeclareFontShape{OT1}{XCharter(0)}{m}{it}{<->ssub * XCharter-TLF/m/it}{}
\DeclareFontShape{OT1}{XCharter(0)}{b}{n}{<->ssub * XCharter-TLF/b/n}{}
\DeclareFontShape{OT1}{XCharter(0)}{bx}{n}{<->ssub * XCharter-TLF/b/n}{}

\usepackage{bm}
\usepackage[table]{xcolor}
\usepackage{xcolor}
\usepackage{algpseudocode}

\usepackage{cancel}
\usepackage{ulem}

\usetikzlibrary{shapes}
\usetikzlibrary{arrows}
\usetikzlibrary {arrows.meta}
\usetikzlibrary{calc,positioning}
\usepgfplotslibrary{fillbetween}
\usetikzlibrary{angles,quotes}
\usetikzlibrary{tikzmark}


\definecolor{PastelRed}{HTML}{FFC1CC}
\definecolor{PastelBlue}{HTML}{C2F0FF}
\definecolor{PastelBlue2}{HTML}{3182BD}

\definecolor{PastelPink}{HTML}{FFCBE1}
\definecolor{PastelGreen}{HTML}{D6E5BD}
\definecolor{PastelYellow}{HTML}{F9E1A8}
\definecolor{PastelBlue}{HTML}{BCD8EC}
\definecolor{PastelPurple}{HTML}{DCCCEC}
\definecolor{PastelOrange}{HTML}{FFDAB4}


\renewcommand{\footnotesize}{\tiny}

\newcommand{\mespor}{
  \ifcase\month
    \or janeiro
    \or fevereiro
    \or março
    \or abril
    \or maio
    \or junho
    \or julho
    \or agosto
    \or setembro
    \or outubro
    \or novembro
    \or dezembro
  \fi
}
% \usepackage{csquotes}
% \usepackage{minted}
% \usemintedstyle{xcode}
\definecolor{vscLightBackground}{HTML}{FFFFFF}
\definecolor{vscLightText}{HTML}{000000}
\definecolor{vscLightGreen}{HTML}{008000}        % For Comments AND Numbers
\definecolor{vscLightRed}{HTML}{A31515}          % For Strings
\definecolor{vscLightTeal}{HTML}{267F99}
\definecolor{vscBlue}{HTML}{0000FF}% For main keywords (if, for, def...)
\definecolor{vscLightPytorch}{HTML}{800080}      % For PyTorch keywords (dark magenta)

\lstdefinestyle{vscode-light}{
    language=Python,
    backgroundcolor=\color{vscLightBackground},
    basicstyle=\ttfamily\scriptsize\color{vscLightText},
    keywordstyle=\color{vscLightTeal},
    commentstyle=\color{vscLightGreen},
    stringstyle=\color{vscLightRed},
    % Style for PyTorch keywords (using keyword group 2)
    keywordstyle=[2]\color{vscLightPytorch}, 
    morekeywords={self, True, False, None},
    morekeywords=[2]{
        torch, Tensor, device, dtype, to, nn, optim, autograd, utils, data,
        Module, Linear, Conv2d, ReLU, CrossEntropyLoss, Adam, SGD,
        Dataset, DataLoader, randn, zeros, ones, arange, cat, stack,
        view, reshape, sum, mean, max, min, matmul, no_grad, backward
    },
    keywordstyle=[3]\color{vscBlue}, 
    morekeywords=[3]{
        model, loss, y, y_hat, x , optimizer
    },
    frame=single,
    tabsize=1,
    showstringspaces=false,
    breaklines=true,
    captionpos=b,
    extendedchars=true,
    numbers=left,
}

\usepackage[cache=true]{minted}
% minted v3+ (TeX Live 2024+) auto-detects the output directory.
% minted v2 (Ubuntu apt TeX Live 2023) needs it set explicitly to
% match latexmkrc's $out_dir = 'build'.
\makeatletter
\@ifpackagelater{minted}{2024/01/01}{}{%
  \def\minted@outputdir{build/}%
}
\makeatother
\setminted{
    fontsize=\scriptsize,
    style=vs,
    breaklines=true,
    numbers=left,
    numbersep=5pt,
    frame=single,
}